\chapter*{APPENDICES}
\addcontentsline{toc}{chapter}{APPENDICES}

{\baselineskip=1\baselineskip
	\centering
	
	\textbf{APPENDIX A}\\[1em]
	\textbf{GANTT CHART}
	\begin{figure}[H]
		\centering
		\makebox[\textwidth][c]{%
			\includegraphics[width=1.4\textwidth]{figures/Gantt-Chart.pdf}%
		}
		% \caption{Gantt Chart}
		% \label{fig:GanttChart}
	\end{figure}
	\newpage

	
	\textbf{APPENDIX B}\\[1em]
	\textbf{FUNCTIONAL TESTING FORM}
	
	\begin{table}[H]
		\centering
		\renewcommand{\arraystretch}{1.3}
		\caption*{\textbf{A. Hardware and Data Capture}}
		\resizebox{\textwidth}{!}{%
			\begin{tabular}{|p{1.2cm}|p{4.2cm}|p{4.2cm}|p{2cm}|p{2cm}|}
				\hline
				\textbf{Test Case No.} & \textbf{Test Description} & \textbf{Expected Result} &
				\multicolumn{2}{|c|}{\textbf{Does it achieve the expected result?}} \\ \cline{4-5}
				& & & \textbf{Yes} & \textbf{No} \\ \hline
				1 & Sensor powers on and initializes correctly. & Sensor activates and displays readiness. & & \\ \hline
				2 & The sensor captures soil parameters (N, P, K, pH, salinity, temperature, moisture). & Accurate readings displayed on serial monitor or app. & & \\ \hline
				3 & Multiple readings under the same conditions are consistent. & Readings show minimal variance (±5\%). & & \\ \hline
				4 & ESP32 collects and processes sensor data properly. & Data packets successfully generated. & & \\ \hline
				5 & The LoRa module transmits data to the gateway. & Data successfully received at gateway. & & \\ \hline
				6 & Data stored correctly in the Firebase database. & Database shows complete and valid entries. & & \\ \hline
			\end{tabular}
		}
	\end{table}
	
	\begin{table}[H]
		\centering
		\renewcommand{\arraystretch}{1.3}
		\caption*{\textbf{B. Mobile Application (Farmer Side)}}
		\resizebox{\textwidth}{!}{%
			\begin{tabular}{|c|p{4.5cm}|p{4.5cm}|c|c|}
				\hline
				\textbf{Test Case No.} & \textbf{Test Description} & \textbf{Expected Result} &
				\multicolumn{2}{|c|}{\textbf{Does it achieve the expected result?}} \\ \cline{4-5}
				& & & \textbf{Yes} & \textbf{No} \\ \hline
				1 & User login and authentication. & Users can log in using valid credentials. & & \\ \hline
				2 & Soil health readings display correctly. & Dashboard shows latest sensor readings. & & \\ \hline
				3 & Crop recommendation generation. & App suggests appropriate crops based on soil data. & & \\ \hline
				4 & Fertilizer recommendation with dosage and schedule. & Correct fertilizer type and amount displayed. & & \\ \hline
				5 & Multiple crop recommendations work without error. & The system handles multiple queries smoothly. & & \\ \hline
				6 & Offline mode and data synchronization. & Data syncs automatically once reconnected. & & \\ \hline
				7 & Farm information update. & Farmers can add/edit farm details successfully. & & \\ \hline
			\end{tabular}
		}
	\end{table}
	
	\begin{table}[H]
		\centering
		\renewcommand{\arraystretch}{1.3}
		\caption*{\textbf{C. Web Application (DA / Administrator Side)}}
		\resizebox{\textwidth}{!}{%
			\begin{tabular}{|c|p{4.5cm}|p{4.5cm}|c|c|}
				\hline
				\textbf{Test Case No.} & \textbf{Test Description} & \textbf{Expected Result} &
				\multicolumn{2}{|c|}{\textbf{Does it achieve the expected result?}} \\ \cline{4-5}
				& & & \textbf{Yes} & \textbf{No} \\ \hline
				1 & Admin login and authentication. & DA officials can securely log in. & & \\ \hline
				2 & Real-time soil health dashboard. & Data from sensors updates in real time. & & \\ \hline
				3 & GIS mapping visualization. & Farm locations shown accurately on the map. & & \\ \hline
				4 & Farmer profile management. & Officials can view and edit farmer data. & & \\ \hline
				5 & Crop rules and recommendation management. & Admins can modify rules and save updates. & & \\ \hline
				6 & Web responsiveness. & The system is accessible and functional on all devices. & & \\ \hline
			\end{tabular}
		}
	\end{table}
	
	\begin{table}[H]
		\centering
		\renewcommand{\arraystretch}{1.3}
		\caption*{\textbf{D. System-Wide Features}}
		\resizebox{\textwidth}{!}{%
			\begin{tabular}{|c|p{4.5cm}|p{4.5cm}|c|c|}
				\hline
				\textbf{Test Case No.} & \textbf{Test Description} & \textbf{Expected Result} &
				\multicolumn{2}{|c|}{\textbf{Does it achieve the expected result?}} \\ \cline{4-5}
				& & & \textbf{Yes} & \textbf{No} \\ \hline
				1 & End-to-end data flow. & Data moves seamlessly from sensor → gateway → Firebase → apps. & & \\ \hline
				2 & Error handling for missing/corrupt data. & System displays proper error messages or retries. & & \\ \hline
				3 & Security and access control. & Only authorized users can access restricted features. & & \\ \hline
				4 & Multi-user capability. & Multiple users can operate simultaneously without system crash. & & \\ \hline
			\end{tabular}
		}
	\end{table}
	
	\newpage
	
	\textbf{APPENDIX C}\\[1em]
	\textbf{USABILITY TESTING SCALE}
	
	\begin{table}[H]
		\centering
		\renewcommand{\arraystretch}{1.3}
		\resizebox{\textwidth}{!}{%
			\begin{tabular}{|c|p{8cm}|c|c|c|c|c|}
				\hline
				\multicolumn{7}{|c|}{\textbf{System Usability Scale (Farmers)}} \\ \hline
				\textbf{No.} & \centering\textbf{Item} & \multicolumn{5}{c|}{\textbf{Frequency}} \\ \hline
				&  & \textbf{SD} & \textbf{D} & \textbf{N} & \textbf{A} & \textbf{SA} \\ \hline
				1 & I think that I would like to use the TANIM mobile app frequently for my farming activities. & & & & & \\ \hline
				2 & I found the TANIM mobile app unnecessarily complex to use. & & & & & \\ \hline
				3 & I thought the app was easy to use for checking my soil's NPK and pH levels. & & & & & \\ \hline
				4 & I think that I would need the support of a technical person to be able to understand the soil parameter and the crop recommendations. & & & & & \\ \hline
				5 & I found the various functions (weather, soil health, and fertilizer schedules) were well integrated. & & & & & \\ \hline
				6 & I thought there was too much inconsistency in how the soil health data (NPK, pH, moisture, temperature, and salinity) and recommendations were displayed. & & & & & \\ \hline
				7 & I would imagine that most farmer would learn to use this app very quickly. & & & & & \\ \hline
				8 & I found the process of viewing my farm very complex. & & & & & \\ \hline
				9 & I felt very confident using the app to make decisions about which crop and fertilizer to apply. & & & & & \\ \hline
				10 & I needed to learn a lot of things before I can use this system. & & & & & \\ \hline
			\end{tabular}
		}
	\end{table}
	
		\begin{table}[H]
		\centering
		\renewcommand{\arraystretch}{1.3}
		\resizebox{\textwidth}{!}{%
			\begin{tabular}{|c|p{8cm}|c|c|c|c|c|}
				\hline
				\multicolumn{7}{|c|}{\textbf{System Usability Scale (Farmers)}} \\ \hline
				\textbf{No.} & \centering\textbf{Item} & \multicolumn{5}{c|}{\textbf{Frequency}} \\ \hline
				&  & \textbf{SD} & \textbf{D} & \textbf{N} & \textbf{A} & \textbf{SA} \\ \hline
				1 & I think that I would like to use the TANIM mobile app frequently for my farming activities. & & & & & \\ \hline
				2 & I found the TANIM mobile app unnecessarily complex to use. & & & & & \\ \hline
				3 & I thought the app was easy to use for checking my soil's NPK and pH levels. & & & & & \\ \hline
				4 & I think that I would need the support of a technical person to be able to understand the soil parameter and the crop recommendations. & & & & & \\ \hline
				5 & I found the various functions (weather, soil health, and fertilizer schedules) were well integrated. & & & & & \\ \hline
				6 & I thought there was too much inconsistency in how the soil health data (NPK, pH, moisture, temperature, and salinity) and recommendations were displayed. & & & & & \\ \hline
				7 & I would imagine that most farmer would learn to use this app very quickly. & & & & & \\ \hline
				8 & I found the process of viewing my farm very complex. & & & & & \\ \hline
				9 & I felt very confident using the app to make decisions about which crop and fertilizer to apply. & & & & & \\ \hline
				10 & I needed to learn a lot of things before I can use this system. & & & & & \\ \hline
			\end{tabular}
		}
	\end{table}
	
	\newpage
	
	\textbf{APPENDIX D}\\[1em]
	\textbf{TESTING CONDUCTED}
	\begin{figure}[H]
		\centering
		\includegraphics[width=0.9\textwidth]{figures/Testing.pdf}
	\end{figure}
	
	\begin{figure}[H]
		\centering
		\includegraphics[width=0.9\textwidth]{figures/Testing1.pdf}
	\end{figure}
	
	\begin{figure}[H]
		\centering
		\includegraphics[width=0.9\textwidth]{figures/Testing2.pdf}
	\end{figure}
	
	\begin{figure}[H]
		\centering
		\includegraphics[width=0.9\textwidth]{figures/Testing3.pdf}
	\end{figure}
	
	\begin{figure}[H]
		\centering
		\includegraphics[width=0.9\textwidth]{figures/Testing4.pdf}
	\end{figure}
	
	
	\newpage
	
	\textbf{APPENDIX E}\\[1em]
	\textbf{LETTER OF CONSENT FOR INTERVIEW}
	
	\begin{figure}[H]
		\centering
		\includegraphics[width=0.9\textwidth]{figures/RequestInterview.pdf}
	\end{figure}
	
	\newpage
	
	\textbf{APPENDIX F}\\[1em]
	\textbf{LETTER OF REQUEST FOR SOIL TESTING KIT}
	\begin{figure}[H]
		\centering
		\includegraphics[width=0.9\textwidth]{figures/RequestLetterDA.pdf}
	\end{figure}
	
	\newpage
	
	\textbf{APPENDIX G}\\[1em]
	\textbf{DOCUMENTS REVIEWED}
	\begin{figure}[H]
		\includegraphics[width=0.9\textwidth]{figures/Documents_Reviewed1.pdf}
	\end{figure}
	
	\begin{figure}[H]
		\includegraphics[width=1\textwidth]{figures/Documents_Reviewed2.pdf}
	\end{figure}

	\begin{figure}[H]
		\includegraphics[width=1\textwidth]{figures/Documents_Reviewed3.pdf}
	\end{figure}

	\begin{figure}[H]
		\includegraphics[width=1\textwidth]{figures/Documents_Reviewed4.pdf}
	\end{figure}
	
	\begin{figure}[H]
		\includegraphics[width=1\textwidth]{figures/Documents_Reviewed5.pdf}
	\end{figure}
	
	\begin{figure}[H]
		\includegraphics[width=1\textwidth]{figures/Documents_Reviewed6.pdf}
	\end{figure}
	
	\newpage
	
	\textbf{APPENDIX H}\\[1em]
	\textbf{INTERVIEW QUESTIONS}
	\begin{longtable}{|p{3cm}|p{5cm}|p{7cm}|}
	\caption{Interview with DA Senior Agriculturist}
	\label{tab:da_interview_full} \\
	
	\hline
	\textbf{Category} & \textbf{Item} & \textbf{Details} \\ \hline
	\endfirsthead
	
	\hline
	\textbf{Category} & \textbf{Item} & \textbf{Details} \\ \hline
	\endhead
	
	\hline
	\multicolumn{3}{|r|}{\textit{Continued on next page}} \\ \hline
	\endfoot
	
	\hline
	\endlastfoot
	
	\multicolumn{3}{|c|}{Interview Information} \\ \hline
	
	Info & Interviewers & Apuya, Christian James E.; Bantaculo, Ray Simon; Diaz, Albert Jhon; Guardaquivil, Alex Jr; Salvo, Robert Roy \\ \hline
	Info & Location & Department of Agriculture, Regional Field Office 10 \\ \hline
	Info & Date & January 27, 2026; February 23, 2026; March 13, 2026 \\ \hline
	Info & Name & Karl Roman Wong \\ \hline
	Info & Address & DA-10 \\ \hline
	Info & Sex & Male \\ \hline
	Info & Occupation & DA Senior Agriculturist \\ \hline
	
	\multicolumn{3}{|c|}{Interview Questions and Responses} \\ \hline
	
	Q1 & How do farmers know the soil nutrients in their agricultural lands? & Ang pag test sa yuta naa tay duha ka methods which is ang Laboratory Results Testing and the Soil Test Kits (STK) for testing, and the output lang ga differ. Sa Laboratory is naa gyud didto ang quantitative nga result sa pH, and NPK, then sa STK purely qualitative to based from the color nga result which is LOW, MEDIUM, HIGH. \\ \hline
	
	Q2 & If ever magpa test ang farmer, unsa ilang method for both ways? & Pag STK lang ang gusto sa farmer, need lang nila ug agreement with the DA to borrow the STK. Meanwhile, sa laboratory results, they will need to transport or bring ilang soil sample diri sa atung laboratory para ma test ilang soil pero naa lang payment since ga involve na siya ug instruments. \\ \hline
	
	Q3 & What is the proper method in order to test a soil sample? & Per hectare mana sila gapa test then need lang nila mag kalut ug atleast 3 ka bangag for atleast 12 inches or 1 foot, and didto kuha sila ug soil sample then mix it with the other soil sample from other holes. Pero ga depende ang ka lalom sa crops nga itanom nila kay naa man tay classification sa crops which is shallow rooted and deep rooted nga mga crops or plants. Pero sa part ninyo focus raman mo sa mga shallow rooted. Also, naa pud tay classification if lowland or highland ang lugar, but since you mentioned nga bukidnon mo so most of it is highland lang siya. \\ \hline
	
	Q5 & How often should our system monitor soil conditions to provide reliable crop recommendations? & Kung laboratory testing, okay ra siya for 1--2 years, pero kung gusto gyud updated ang data, mas maayo every cropping cycle, which is pre and post planting. \\ \hline
	
	Q6 & What type of data should be collected to improve its crop recommendation accuracy? & Kinahanglan naay baseline data sa soil ug crops. Pareha sa among ginagamit nga apps, ang data gikan sa field trials ug techno demo. Importante pud ang historical data para mas accurate ang recommendations. \\ \hline
	
	Q7 & Which crops should TANIM prioritize when recommending based on soil conditions? & Sa among region, priority ang rice ug corn ug cassave, pero dapat various crops maapil kay lain-lain man silag nutrient requirements. \\ \hline
	
	Q8 & What challenges might be face in terms of farmer adoption of soil monitoring systems? & Bisan naa nay mga apps karon, gamay ra ang nag adopt. Lisod pud usahay kay ang farmers naanad na sa ilang sariling pamaagi. Kinahanglan ang system kay simple ug dali masabtan. \\ \hline
	
	Q9 & How do farmers currently decide on fertilizer application, and how can Fertright improve this process? & Kasagaran sa farmers kay experience-based ra ilang pag apply ug fertilizer. Usahay dili sila musunod sa recommendations kay naanad na sila sa ilang style. Makatabang ang Fertright kung makahatag siya ug klaro, sakto, ug masaligan nga guide. \\ \hline
	
	Q10 & How can Fertright ensure that its fertilizer recommendations are practical for farmers? & Dapat ang recommendations kay precise ug accurate. Dili dapat komplikado kay ang farmers ganahan ug simple nga instructions nga dali masunod. \\ \hline
	
	Q11 & What features should be included to be usable for farmers in rural areas? & Importante nga naa siyay offline functionality kay daghan farmers walay internet. Dapat pud user-friendly ug dili lisod gamiton. \\ \hline
	
	Q12 & What improvements in fertilizer management to promote sustainable farming? & Pwede siya makatabang sa balanced fertilization, paggamit ug organic fertilizers, ug soil conservation practices. Importante ni para magdugay ang kaayo sa yuta ug mudako ang ani. \\ \hline
	
\end{longtable}
	
	\newpage
	\textbf{APPENDIX I}\\[1em]
	\textbf{CALIBRATION OF SOIL SENSOR USING SOIL TESTING KIT}
	\label{app:stk_results}
	\begin{figure}[H]
		\centering
		\includegraphics[width=0.85\textwidth]{figures/SoilTestKit1.pdf}
	\end{figure}
	
	\begin{figure}[H]
		\centering
		\includegraphics[width=0.85\textwidth]{figures/SoilTestKit2.pdf}
	\end{figure}
	
	\newpage
	
	\textbf{APPENDIX J}\\[1em]
	\textbf{DOCUMENTATION OF THE IMPLEMENTATION}
	
	\begin{figure}[H]
		\centering
		\includegraphics[width=0.8\textwidth]{figures/ImplementationDocumentation.pdf}
		\label{fig:implementation_documentation}
	\end{figure}
	\newpage
	
	\textbf{APPENDIX K}\\[1em]
	\textbf{FUNCTIONAL AND USABILITY TESTING CONDUCTED}
	\begin{figure}[H]
		\centering
		\includegraphics[width=0.85\textwidth]{figures/SUS1.pdf}
	\end{figure}
	
	\begin{figure}[H]
		\centering
		\includegraphics[width=0.85\textwidth]{figures/SUS2.pdf}
	\end{figure}
	
	\begin{figure}[H]
		\centering
		\includegraphics[width=0.85\textwidth]{figures/SUS3.pdf}
	\end{figure}
	
	\begin{figure}[H]
		\centering
		\includegraphics[width=0.85\textwidth]{figures/SUS4.pdf}
	\end{figure}
	
	\begin{figure}[H]
		\centering
		\includegraphics[width=0.85\textwidth]{figures/SUS5.pdf}
	\end{figure}
	
	\begin{figure}[H]
		\centering
		\includegraphics[width=0.85\textwidth]{figures/SUS6.pdf}
	\end{figure}
	
	\begin{figure}[H]
		\centering
		\includegraphics[width=0.85\textwidth]{figures/SUS7.pdf}
	\end{figure}
	
	\begin{figure}[H]
		\centering
		\includegraphics[width=0.85\textwidth]{figures/SUS8.pdf}
	\end{figure}
	
	\begin{figure}[H]
		\centering
		\includegraphics[width=0.85\textwidth]{figures/SUS9.pdf}
	\end{figure}
	
	\begin{figure}[H]
		\centering
		\includegraphics[width=0.85\textwidth]{figures/SUS10.pdf}
	\end{figure}
	
	\newpage

	
	\textbf{APPENDIX L}\\[1em]
	\textbf{SENSOR DATA SHEET}
	\begin{figure}[H]
		\centering
		\includegraphics[width=0.85\textwidth]{figures/DataSheet.pdf}
	\end{figure}
	
	\begin{figure}[H]
		\centering
		\includegraphics[width=0.85\textwidth]{figures/DataSheet2.pdf}
	\end{figure}
	
	\newpage
	
	\textbf{APPENDIX M}\\[1em]
	\textbf{TURNITIN RESULTS}
	
	\newpage
	
	\textbf{APPENDIX N}\\[1em]
	\textbf{GRAMMARIAN RESULTS}
	
}
